\documentclass[reqno,11pt]{amsart}
\usepackage{math327-lecture}
\usepackage{needspace}
\lectureinfo{10}{Sep. 25, 2026}{Cosets and Equivalence Relations}

% Based on handwritten Lecture 8, Sep. 19, 2025.
\begin{document}
\makelecturetitle


\section{Left and right cosets}

\begin{define}[Cosets and representatives]
Let $H\leq G$ and $g\in G$. The \emph{left coset} and \emph{right coset}
of $H$ represented by $g$ are
\begin{equation}
  gH=\{gh:h\in H\},\qquad Hg=\{hg:h\in H\}.
\end{equation}
We call $g$ a \emph{coset representative}. The set of all left cosets is
\begin{equation}
  G/H=\{gH:g\in G\}.
\end{equation}
\end{define}

A coset is a subset of $G$; its representative is an element of $G$.
Different representatives can give the same coset. If $g\in H$, then
$gH=H$; we prove this in Lemma~\ref{lem:coset-equality} below, and the
same argument with multiplication on the right shows $Hg=H$.

Cosets give a new way to express normality. For a subgroup $N\leq G$ and
$g\in G$,
\begin{equation}
	gNg^{-1}=N\quad\Longleftrightarrow\quad gN=Ng.
\end{equation}
Why? Assume the right-hand side. Multiplying $gN=Ng$ on the right by
$g^{-1}$ gives $gNg^{-1}=Ngg^{-1}=N$. Conversely, if $gNg^{-1}=N$, then
multiplying on the right by $g$ gives $gN=Ng$. Thus $N$ is normal exactly
when every left coset of $N$ equals the corresponding right coset.

In the examples below we write
\[
  S_3=\{1,x,x^2,y,xy,x^2y\},\qquad
  x^3=1,\quad y^2=1,\quad yx=x^2y.
\]

\begin{example}[Cosets of the subgroup of order three in $S_3$]
Take $N=\{1,x,x^2\}$. For its elements, we obtain
\begin{align*}
  1N&=\{1,x,x^2\}=N,\\
  xN&=\{x,x^2,x^3\}=N,\\
  x^2N&=\{x^2,x^3,x^4\}=N.
\end{align*}
For the other three elements, use $yx=x^2y$, $yx^2=xy$, and $xyx=y$:
\begin{align*}
  yN&=\{y,yx,yx^2\}=\{y,x^2y,xy\},\\
  xyN&=\{xy,xyx,xyx^2\}=\{xy,y,x^2y\}=yN,\\
  x^2yN&=\{x^2y,x^2yx,x^2yx^2\}=\{x^2y,xy,y\}=yN.
\end{align*}
There are exactly two distinct left cosets:
\[
  S_3/N=\{N,yN\},\qquad S_3=N\sqcup yN.
\]
The symbol $\sqcup$ denotes a disjoint union. Since $N$ is normal, the
right coset represented by any $g$ equals the corresponding left coset.
\end{example}

\begin{example}[Cosets of a subgroup of order two in $S_3$]
Now take the non-normal subgroup $H=\{1,y\}$. The left cosets are
\begin{align*}
  H=yH&=\{1,y\},\\
  xH=xyH&=\{x,xy\},\\
  x^2H=x^2yH&=\{x^2,x^2y\}.
\end{align*}
Consequently,
\[
  S_3/H=\{H,xH,x^2H\},\qquad
  S_3=H\sqcup xH\sqcup x^2H.
\]
Here left and right cosets differ:
\[
  xH=\{x,xy\},\qquad Hx=\{x,yx\}=\{x,x^2y\}.
\]
Since $xy\neq x^2y$, we have $xH\neq Hx$, consistent with the criterion
above: $H$ is not normal. Nevertheless, the left cosets still partition
$S_3$, as we now prove in general.
\end{example}

\section{Cosets form a partition}

\begin{define}[Partition]
A \emph{partition} of a set $X$ is a collection of nonempty subsets
$\{X_\alpha\}_{\alpha\in I}$ such that
\[
  X=\bigcup_{\alpha\in I}X_\alpha,
  \qquad X_\alpha\cap X_\beta=\varnothing\quad(\alpha\neq\beta).
\]
We write $X=\bigsqcup_{\alpha\in I}X_\alpha$.
\end{define}

\begin{lemma}[Equality of cosets]\label{lem:coset-equality}
Let $H\leq G$ and $g_1,g_2\in G$. The following are equivalent:
\begin{enumerate}[label=\textbf{(\arabic*)}]
  \item $g_1H\cap g_2H\neq\varnothing$;
  \item $g_1H=g_2H$;
  \item $g_2^{-1}g_1\in H$;
  \item $g_1^{-1}g_2\in H$.
\end{enumerate}
\end{lemma}
\begin{proof}
First consider the special case $g_2=e_G$. We show that
\[
  gH\cap H\neq\varnothing
  \quad\Longleftrightarrow\quad g\in H
  \quad\Longleftrightarrow\quad gH=H.
\]
We prove $g\in H\Rightarrow gH=H\Rightarrow gH\cap H\neq\varnothing
\Rightarrow g\in H$.
\begin{itemize}
  \item If $g\in H$, closure gives $gH\subseteq H$. Conversely, every
  $h\in H$ can be written $h=g(g^{-1}h)$, with $g^{-1}h\in H$, so
  $H\subseteq gH$. Thus $gH=H$.
  \item If $gH=H$, then $gH\cap H=H$, which is nonempty since
  $e_G\in H$.
  \item If $z\in gH\cap H$, write $z=gh$ with $h\in H$. Then
  $g=zh^{-1}\in H$.
\end{itemize}
This completes the special case.

For general $g_1,g_2$, left multiplication by $g_2^{-1}$ is a bijection
of $G$. Hence
\begin{align*}
  g_1H\cap g_2H\neq\varnothing
  &\quad\Longleftrightarrow\quad
  (g_2^{-1}g_1)H\cap H\neq\varnothing\\
  &\quad\Longleftrightarrow\quad g_2^{-1}g_1\in H\\
  &\quad\Longleftrightarrow\quad (g_2^{-1}g_1)H=H\\
  &\quad\Longleftrightarrow\quad g_1H=g_2H.
\end{align*}
Finally, $g_2^{-1}g_1\in H$ if and only if its inverse
$g_1^{-1}g_2$ belongs to $H$.
\end{proof}

\begin{proposition}
The distinct left cosets of $H$ form a partition of $G$.
\end{proposition}
\begin{proof}
Every $g\in G$ belongs to $gH$, since $g=ge_G$ and $e_G\in H$.
Thus the cosets are nonempty and cover $G$. By
Lemma~\ref{lem:coset-equality}, two cosets that intersect are equal, so
distinct cosets are disjoint.
\end{proof}

\begin{warningbox}
$G/H$ denotes the \emph{set of left cosets}, not a set of representatives.
If $R\subseteq G$ contains exactly one representative of each left coset,
the precise formulas are
\begin{equation}
  G/H=\{rH:r\in R\},\qquad G=\bigsqcup_{r\in R}rH.
\end{equation}
The representative set $R$ is not unique. For $N=\{1,x,x^2\}$ in $S_3$,
both $\{1,y\}$ and $\{x,xy\}$ are representative sets; for $H=\{1,y\}$,
both $\{1,x,x^2\}$ and $\{y,xy,x^2y\}$ are. When $H$ is normal, we will
later give $G/H$ a group operation.
\end{warningbox}

\section{Example: cosets of \texorpdfstring{$8\Z$}{8Z}}\label{sec:8Z}

In the additive group $G=\Z$, take $H=8\Z$. Cosets are written additively:
\begin{equation}
  a+8\Z=\{a+8k:k\in\Z\}.
\end{equation}
For instance,
\begin{align*}
  0+8\Z&=\{\ldots,-16,-8,0,8,16,\ldots\},\\
  1+8\Z&=\{\ldots,-15,-7,1,9,17,\ldots\},\\
  7+8\Z&=\{\ldots,-9,-1,7,15,23,\ldots\}.
\end{align*}
Arranging the integers by their remainders gives eight disjoint columns:
\[
\begin{array}{c|rrrrrrrr}
  \text{remainder}&0&1&2&3&4&5&6&7\\ \hline
  &\vdots&\vdots&\vdots&\vdots&\vdots&\vdots&\vdots&\vdots\\
  &16&17&18&19&20&21&22&23\\
  &8&9&10&11&12&13&14&15\\
  &0&1&2&3&4&5&6&7\\
  &-8&-7&-6&-5&-4&-3&-2&-1\\
  &-16&-15&-14&-13&-12&-11&-10&-9\\
  &\vdots&\vdots&\vdots&\vdots&\vdots&\vdots&\vdots&\vdots
\end{array}
\]
Each column is one coset. The equality criterion becomes
\begin{equation}
  a+8\Z=b+8\Z
  \quad\Longleftrightarrow\quad a-b\in8\Z
  \quad\Longleftrightarrow\quad a\equiv b\pmod8.
\end{equation}
For example, $0+8\Z=8+8\Z$ and $1+8\Z=17+8\Z$.
Every integer has exactly one remainder in $\{0,1,\ldots,7\}$, so
\[
  \Z=\bigsqcup_{r=0}^7(r+8\Z),\qquad
  \Z/8\Z=\{0+8\Z,1+8\Z,\ldots,7+8\Z\}.
\]
Writing $\overline a=a+8\Z$, we may abbreviate this as
\[
  \Z/8\Z=\{\overline0,\overline1,\ldots,\overline7\}.
\]
Thus an infinite group can have a finite set of cosets.

%
%\section{Equivalence relations and partitions}
%
%An equivalence relation is a way to declare elements of a set equivalent
%while retaining the basic properties of equality.
%
%\begin{define}[Equivalence relation]
%Let $X$ be a set. A relation $\sim$ on $X$ is an \emph{equivalence relation}
%if, for all $a,b,c\in X$, it satisfies:
%\begin{enumerate}[label=\textbf{(\arabic*)}]
%  \item \emph{Reflexivity:} $a\sim a$.
%  \item \emph{Symmetry:} if $a\sim b$, then $b\sim a$.
%  \item \emph{Transitivity:} if $a\sim b$ and $b\sim c$, then $a\sim c$.
%\end{enumerate}
%The \emph{equivalence class} of $a\in X$ is
%\[
%  C_a=\{z\in X:z\sim a\}.
%\]
%\end{define}
%
%\begin{proposition}
%The distinct equivalence classes form a partition of $X$. More precisely,
%\[
%  C_a\cap C_b\neq\varnothing
%  \quad\Longleftrightarrow\quad C_a=C_b
%  \quad\Longleftrightarrow\quad a\sim b.
%\]
%\end{proposition}
%\begin{proof}
%Reflexivity gives $a\in C_a$, so the classes are nonempty and cover $X$.
%If $z\in C_a\cap C_b$, then $z\sim a$ and $z\sim b$. Symmetry and
%transitivity give $a\sim b$. If $a\sim b$ and $w\in C_a$, then
%$w\sim a\sim b$, so $w\in C_b$. This gives $C_a\subseteq C_b$; symmetry
%gives the reverse inclusion. Finally, if $C_a=C_b$, their intersection
%contains $a$. Thus classes are either equal or disjoint.
%\end{proof}
%
%Choose a set $R\subseteq X$ containing exactly one representative of each
%class. Then
%\[
%  X=\bigsqcup_{a\in R}C_a.
%\]
%Conversely, any partition of $X$ defines an equivalence relation: declare
%$a\sim b$ when $a$ and $b$ lie in the same part.
%
%\begin{example}[Congruence modulo eight]
%On $\Z$, define $a\sim b$ if $8\mid(b-a)$. This is an equivalence
%relation (transitivity uses $c-a=(c-b)+(b-a)$), and its class
%$C_a=a+8\Z$ is exactly the column containing $a$ in Section~\ref{sec:8Z}.
%Besides $R=\{0,1,\ldots,7\}$ we could use $R'=\{-3,-2,\ldots,4\}$:
%\[
%  \Z=\bigsqcup_{a\in R}C_a=\bigsqcup_{a\in R'}C_a.
%\]
%The classes are fixed by the relation; the representatives are not.
%\end{example}
%
%\section{Cosets as equivalence classes}
%
%Let $H\leq G$. Define a relation on $G$ by
%\begin{equation}
%  a\sim b\quad\Longleftrightarrow\quad b^{-1}a\in H.
%\end{equation}
%It is an equivalence relation:
%\begin{enumerate}[label=\textbf{(\arabic*)}]
%  \item $a^{-1}a=e_G\in H$, so $a\sim a$.
%  \item If $b^{-1}a\in H$, then $(b^{-1}a)^{-1}=a^{-1}b\in H$, so
%  $b\sim a$.
%  \item If $b^{-1}a\in H$ and $c^{-1}b\in H$, then
%  \[
%    c^{-1}a=(c^{-1}b)(b^{-1}a)\in H,
%  \]
%  so $a\sim c$.
%\end{enumerate}
%The class of $a$ is precisely its left coset:
%\begin{equation}
%  C_a=\{g\in G:a^{-1}g\in H\}=aH.
%\end{equation}
%Indeed, $a^{-1}g=h\in H$ if and only if $g=ah$.
%This gives another explanation of why left cosets partition $G$.
%
%\begin{remark}[The relation for right cosets]
%To obtain right cosets, use the relation
%\[
%  a\sim_R b\quad\Longleftrightarrow\quad ab^{-1}\in H.
%\]
%Then the class of $a$ is $\{g\in G:ga^{-1}\in H\}=Ha$.
%By contrast, replacing $b^{-1}a\in H$ with $a^{-1}b\in H$ does
%\emph{not} change the left-coset relation: these two elements are inverses,
%and $H$ is closed under inverses.
%\end{remark}

%
%\section{Lagrange's theorem}
%
%\begin{lemma}
%For every subgroup $H\leq G$ and every $g\in G$, the map
%\[
%  H\longrightarrow gH,\qquad h\longmapsto gh,
%\]
%is a bijection.
%\end{lemma}
%\begin{proof}
%It is surjective by the definition of $gH$. If $gh_1=gh_2$, cancellation
%gives $h_1=h_2$, so it is injective. Its inverse sends $z\in gH$ to
%$g^{-1}z\in H$.
%\end{proof}
%
%More generally, left multiplication by $g$ is a bijection $G\to G$.
%In particular, if $H$ is finite, every coset has the same number of
%elements as $H$.
%
%\begin{define}[Index]
%The \emph{index} of a subgroup $H$ in $G$, denoted $[G:H]$, is the number
%of distinct left cosets:
%\[
%  [G:H]=|G/H|.
%\]
%The index may be infinite.
%\end{define}
%
%\begin{theorem}[Lagrange]
%If $G$ is finite and $H\leq G$, then
%\begin{equation}
%  |G|=[G:H]\,|H|.
%\end{equation}
%In particular, $|H|$ divides $|G|$.
%\end{theorem}
%\begin{proof}
%Choose one representative of each left coset, say $g_1,\ldots,g_m$, where
%$m=[G:H]$. Since the cosets partition $G$ and each has $|H|$ elements,
%\[
%  G=\bigsqcup_{i=1}^m g_iH,
%  \qquad |G|=\sum_{i=1}^m|g_iH|=m|H|.
%\]
%\end{proof}
%
%\begin{corollary}
%If $G$ is finite and $a\in G$, then the order of $a$ divides $|G|$.
%\end{corollary}
%\begin{proof}
%Apply Lagrange's theorem to the cyclic subgroup $\langle a\rangle$, whose
%cardinality is the order of $a$.
%\end{proof}
%
%\begin{corollary}
%Every group of prime order $p$ is cyclic, hence isomorphic to the cyclic
%group $C_p$ of order $p$.
%\end{corollary}
%\begin{proof}
%Choose $a\neq e_G$. Its order divides $p$, so it is either $1$ or $p$.
%It cannot be $1$, and hence $|\langle a\rangle|=p=|G|$. Therefore
%$G=\langle a\rangle$.
%\end{proof}
%
%\Needspace{12\baselineskip}
%\section{Towards classifying groups of order six}
%
%Let $|G|=6$. Lagrange's theorem shows that an element can have order
%$1$, $2$, $3$, or $6$. We first establish that elements of orders $2$ and
%$3$ must occur.
%
%\begin{lemma}
%Every finite group of even order contains an element of order $2$.
%\end{lemma}
%\begin{proof}
%Pair each element $a$ with its inverse $a^{-1}$. The distinct sets
%$\{a,a^{-1}\}$ partition $G$ into parts of size one or two. A part has
%size one exactly when $a=a^{-1}$.
%If the identity were the only such element, then $|G|$ would be one plus
%an even number, hence odd. Thus there is some $a\neq e_G$ with
%$a=a^{-1}$. It satisfies $a^2=e_G$ and has order $2$.
%\end{proof}
%
%\begin{lemma}
%If $a^2=e_G$ for every $a\in G$, then $G$ is abelian.
%\end{lemma}
%\begin{proof}
%Every element is its own inverse, so for $a,b\in G$,
%\[
%  ba=b^{-1}a^{-1}=(ab)^{-1}=ab.
%\]
%\end{proof}
%
%\begin{proposition}
%A group of order $6$ contains an element of order $3$.
%\end{proposition}
%\begin{proof}
%Suppose first that every nonidentity element has order $2$.
%The preceding lemma makes $G$ abelian. Choose distinct nonidentity
%elements $a,b$. Then
%\[
%  K=\{e_G,a,b,ab\}=\langle a,b\rangle
%\]
%is a subgroup: since $a,b$ commute and $a^2=b^2=e_G$, every product of
%powers of $a,b$ reduces to one of these four elements. They are distinct:
%$ab=e_G$ would imply $a=b$, while $ab=a$ or $ab=b$ would force $b=e_G$
%or $a=e_G$, respectively. Thus $|K|=4$, contradicting Lagrange's theorem
%because $4$ does not divide $6$.
%
%There is therefore an element $a$ whose order is neither $1$ nor $2$.
%Its order is $3$ or $6$. In the first case, take $x=a$.
%In the second case, take $x=a^2$; then $x^3=e_G$ and $x\neq e_G$, so
%$x$ has order $3$.
%\end{proof}
%
%Choose $x$ of order $3$ and $y$ of order $2$, and put
%$H=\langle x\rangle=\{e_G,x,x^2\}$. Since the elements of $H$ have
%orders $1$ or $3$, we have $y\notin H$. The coset criterion gives
%$yH\neq H$. Both cosets have three elements, so
%\begin{equation}
%  G=H\sqcup yH=\{e_G,x,x^2,y,yx,yx^2\}.
%\end{equation}
%In particular, $x$ and $y$ generate $G$.
%Determining how these generators multiply will be the next step in the
%classification.

\end{document}
